% Lösungen Potenz- und Exponentialfunktionen, Wachstum und Zerfall · Wert und Schwellenwert berechnen · Prüfungs-Fokus MSA (zusammenbau v0.6)
\einheitenkopf{Einheit 3 · Exponentialfunktion aufstellen und auswerten}

\erg{1}{a) $1{,}04^5 \approx 1{,}2167$ \quad b) $440$; $484$; $532{,}4$ \quad c) $t = 6$: $2\,500 \cdot 1{,}06^6 \approx 3\,546$}
\erg{2}{a) $N(30) = 64 \cdot 1{,}25^{30} \approx 51\,699$ – die Aussage stimmt. \quad b) $N(20) = 120 \cdot 2^{20} = 125\,829\,120$ – die Aussage stimmt. \quad c) $x = 12$: $36 \cdot 1{,}04^{12} \approx 57{,}6$; Studie 2: $36 \cdot 1{,}5 = 54$ – Studie 1 sagt für $2033$ den höheren Verbrauch voraus. \quad d) $x = 20$: $120 \cdot 1{,}02^{20} \approx 178{,}3$; Studie 2: $120 \cdot 1{,}4 = 168$ – Studie 1 liegt höher. \quad e) $8 \cdot 0{,}8^{24} \approx 0{,}04\,\mathrm{mg}$ \quad f) $80 \cdot 0{,}75^{24} \approx 0{,}08\,\mathrm{mg}$ \quad g) $800 \cdot 0{,}88^4 \approx 480\,\mathrm{hPa}$ – die Behauptung stimmt. \quad h) $1\,000 \cdot 0{,}85^5 \approx 444$ Lux – die Behauptung stimmt. \quad i) $2027$: $654\,840$ ($2024$); $700\,679$ ($2025$); $749\,726$ ($2026$); $802\,207\,€$ ($2027$) \quad j) $2027$: $8\,487$; $8\,742$; $9\,004$; $9\,274$; $9\,552$ Einwohner ($2027$)}
